\section{磁场}

\subsection{带电粒子在电磁场中运动}

\begin{proset}
    如图所示，在\(x\)轴上方空间存在电场强度大小为\(E\)的匀强电场，电场方向在\(xOy\)平面（纸面）内且与\( + x\)的夹角为\,\qty{53}{\degree}斜向右下方；在\(x\)轴下方空间中，存在垂直\(xOy\)平面（纸面）向外的匀强磁场。一比荷为\(k\)的带正电粒子，以速度\(v_0\)从坐标原点\(O\)沿\( + y\)方向射入电场，并通过\(x\)轴上的\(A\)点后再次回到坐标原点\(O\)，不计粒子重力。已知\(\sin\) \,\qty{53}{\degree}\( = 0.8\)，\(\cos\) \,\qty{53}{\degree}\( = 0.6\)，求：
    \begin{enumerate}[(1)]
        \item 粒子从\(O\)运动到\(A\)点的时间；
        \item 粒子通过\(A\)点时的速度大小\(v_1\)；
        \item 匀强磁场的磁感应强度大小\(B\)；
        \item 粒子第\(n\)次从磁场进入电场时的横坐标。
    \end{enumerate}
    \begin{tikzpicture}[font=\scriptsize,>=latex]
        \draw[->,blue] (-1,0)--(4,0)node[below]{\(x\)};
        \draw[->,blue] (0,-3)--(0,2)node[left]{\(y\)};
        \foreach \x in {-.2,.5,...,3.8}
            \draw[->,blue] ($(\x,0)+(127:1.5)$)--++(-53:1.5);
        \foreach \x in{-.4,.1,...,3.8}
            \foreach \y in{-.3,-.8,-1.3}
                \fill (\x,\y) circle(.03);
        \node at(-.2,-.2){\(O\)};
        \fill (1.5,0)node[above]{\(A\)} circle(.05);
        \node[blue,above] at (1.8,1.2){\(E\)};
    \end{tikzpicture}
\end{proset}
\begin{answer}
    （1）\(\frac{5v_0}{2kE}\)（2）\(\frac{\sqrt{13}}{2}v_0 \)（3）\(\frac{16E}{15v_0}\)（4）\(\frac{15n(n - 1)v_0^2}{8kE}(n =1,2,3\cdots)\)
\end{answer}